% Referee response template.
% Structure only; no responses written, because no report exists.
\documentclass[11pt]{article}
\usepackage[margin=1in]{geometry}
\usepackage{amsmath,amssymb,xcolor}
\newcommand{\TODO}[1]{\textbf{[#1]}}
\newcommand{\refpoint}[1]{\medskip\noindent\textbf{Referee:} \emph{#1}\par\medskip\noindent\textbf{Response:} }
\begin{document}
\begin{center}
{\Large Response to the Referee Report}\\[2mm]
\emph{Symmetry control of exceptional structure in the massive scalar
quasinormal spectrum of doubly rotating five-dimensional Myers--Perry black
holes}\\[2mm]
\TODO{manuscript number}
\end{center}

\noindent We thank the referee for their reading of the manuscript. Changes in
the revised version are marked in \textcolor{blue}{blue}. Below we reproduce
each point in italics and respond beneath it.

\section*{Anticipated points}

The following are objections we expect and for which the evidence already
exists in the manuscript or the public repository. They are listed so that the
response can be specific rather than defensive.

\refpoint{The result is negative; how do you know you simply did not look hard
enough, or that your detector could not fire?}
The detector is validated on closed-form problems with known exceptional
structure ($\lambda^2-t$, $\lambda^3-t$, Jordan blocks) \emph{and} on negative
controls where it must stay silent (avoided crossings, semisimple and
block-diagonal degeneracies). See Supplemental Material, Sec.~S4. Separately,
the exclusion is quantitative: three rescaling-invariant diagnostics are bounded
away from zero, so the statement is a bound rather than an absence of sightings.

\refpoint{Exceptional lines are known in Myers--Perry--de Sitter. Is your result
in conflict?}
No. That mechanism relies on the cosmological horizon and the near-Nariai
P\"oschl--Teller limit, neither of which has a $\Lambda=0$ analogue. This is
stated in Sec.~X of the manuscript and in the novelty audit in the repository.

\refpoint{Why not simply use $|dF/d\omega|$ as the exceptional-point diagnostic?}
Because it is not invariant: the spectral condition is defined up to a
nonvanishing analytic factor, so $|dF/d\omega|$ at a root can be rescaled at
will. We exhibit two conditions with identical zeros whose $|dF/d\omega|$ differ
by more than $10^5$. See Sec.~V.

\refpoint{The minimum gap sits on the boundary of your search box.}
Correct, and stated. A probe beyond the box shows the gap falling further at
larger $s$ but rising again at larger $\mu$, so it does not extrapolate to zero;
we make no claim outside the declared domain.

\refpoint{Can the no-go be promoted to the unseparated problem / to a continuum
theorem?}
The no-go already \emph{is} a statement about the unseparated pencil: the
decomposition is over $(m_1,m_2)$ only and uses the torus action, not separation
of variables in $\theta$. It is not a numerical statement and carries no
discretization. Certification of the \emph{numerical} results is a separate
matter and is deliberately limited; see Supplemental Material, Sec.~S5.

\section*{Referee report}

\refpoint{\TODO{referee point 1}}
\TODO{response}

\refpoint{\TODO{referee point 2}}
\TODO{response}

\end{document}
